\begin{longtable}{L{0.23\linewidth}L{0.68\linewidth}}
\caption{非功能性需求与验证方法}\label{tab:nfr}\\\toprule 编号及属性 & 要求与验证依据\\\midrule\endfirsthead
\toprule 编号及属性 & 要求与验证依据\\\midrule\endhead
NFR-01 正确性 & 结构、逐文件字节及摘要一致；增量结果与目标manifest相符。 验证：A-TC-01,02,03,04,24,33\\
NFR-02 元数据正确性 & 权限允许时按声明恢复；跳过属主必须有计数及诊断。 验证：A-TC-12,13\\
NFR-03 失败可见 & 失败返回真实状态，不发布部分内容为成功；无法验证身份则留存诊断。 验证：A-TC-34-42\\
NFR-04 崩溃一致与幂等 & 单文件原子发布，无变化不新建，基线损坏保守重建；多个文件更新由业务流程协调。 验证：A-TC-21,22,26\\
NFR-05 易用性 & CLI退出码区分、GUI忙碌与错误明确、危险操作确认、预览截断可见。 验证：A-TC-05,07,11,27\\
NFR-06 可维护性 & 共享规则/模式/路径语义，核心与Qt分离，源码检查与接口可定位。 验证：E08及接口表\\
NFR-07 本地安全 & 拒绝危险路径、覆盖和无归属依据的删除；文件系统限制见设计章节。 验证：A-TC-06,08,39-42\\
NFR-08 远端安全 & 服务器身份、账户隔离、限速、不降级与数据锁均需检验。 验证：A-TC-28-31\\
NFR-09 可测试与复现 & 记录源码版本、环境、输入、输出、退出码及跳过原因。 验证：E05、E07及E15\\
NFR-10 资源边界 & 流式编码与传输，长度先校验，固定worker；全树索引仍随条目数增长。 验证：脚本边界检查；编解码性能记录见E15\\
NFR-11 构建与安装 & Make 目标可分别构建客户端和服务端，发行包可在目标环境安装和启动。验证：E15 发行包与持续集成检查。\\
\bottomrule\end{longtable}
